% !TeX program = xelatex
% Complete section-mode example for noteformyself 2.0.
\documentclass[sectionlevel=section,status=draft]{noteformyself}

% Section mode prints only this compact author/date line; it has no cover.
\title{A Standalone Section}
\author{Tianle Yang}
\date{\today}
\version{2.0.0}

\LoadNoteTagRegistry{tag-registry.example.tex}
\notesetup{
  book={Tag System Test Book},
  book-tag=0,
  chapter-tag=0SE,
  status=draft
}

\begin{document}

\maketitle

% =============================================================================
% Published section example
% =============================================================================
\section[
  status=published,
  tag=0SE01,
  label=sec:section-preview
]{A published standalone section}

\begin{definition}[
  title={Readable typography},
  status=published,
  tag=0SE0101,
  label=def:section-preview
]
  A template is readable when its visual hierarchy supports the mathematics
  without competing with it.
\end{definition}

The local references are \cref{sec:section-preview} and
\cref{def:section-preview}.

\subsection[
  status=published,
  tag=0SE0102,
  label=subsec:section-preview
]{A tagged subsection}

Section mode uses a centered page number and no header or footer rules.

% =============================================================================
% Complete environment gallery
% =============================================================================
\section[
  status=draft,
  tag=0SE02,
  label=sec:section-gallery
]{A draft environment gallery}

\subsection[tag=0SE0201,label=subsec:section-gallery]{Every statement style}

\begin{slogan}
  Short independently compiled sections should still share the book's voice.
\end{slogan}

\begin{mainthm}[title={Standalone coherence},label=mainthm:section]
  A standalone section can use the same source as its chapter and book builds.
\end{mainthm}

\begin{definition}[label=def:section-draft]
  A chain complex is a sequence \((C_n,d_n)\) with
  \(d_{n-1}\circ d_n=0\).
\end{definition}

\begin{proposition}[title={Cycles modulo boundaries},label=prop:section]
  The quotient \(\ker d_n/\operatorname{im}d_{n+1}\) is an abelian group.
\end{proposition}

\begin{theorem}[label=thm:section]
  A short exact sequence of complexes induces a long exact sequence in
  homology.
\end{theorem}

\begin{lemma}[title={Connecting map},label=lem:section]
  The connecting morphism is independent of the chosen lift.
\end{lemma}

\begin{corollary}[label=cor:section]
  A quasi-isomorphism induces isomorphisms on every homology group.
\end{corollary}

\begin{conjecture}[label=conj:section]
  An analogous statement holds in the intended derived enhancement.
\end{conjecture}

\begin{question}[title={Naturality},label=ques:section]
  How does the connecting map vary with a morphism of short exact sequences?
\end{question}

\begin{example}[tag=0SE0202,label=eg:section]
  The complex concentrated in degree zero recovers an ordinary module. Its
  displayed draft tag does not appear in \cref{eg:section}.
\end{example}

\begin{exercise}[label=ex:section]
  Write down the connecting map and verify that it is a homomorphism.
\end{exercise}

\begin{construction}[title={Mapping cone},label=constr:section]
  For \(f\colon C_\bullet\to D_\bullet\), set
  \(\operatorname{Cone}(f)_n=D_n\oplus C_{n-1}\) with the usual differential.
\end{construction}

\begin{notation}[label=not:section]
  Write \(H_n(C_\bullet)\) for the \(n\)-th homology group.
\end{notation}

\begin{remark}[title={Shared numbering},label=rmk:section]
  Every ordinary theorem-like environment participates in one local sequence.
\end{remark}

\begin{proof}[Proof sketch of \cref{thm:section}]
  Lift a cycle, apply the differential, and descend the resulting boundary.

  \begin{step}[Choose a lift]\label{step:section-lift}
    Select a representative in the middle complex.
  \end{step}

  \begin{claim}[label=claim:section]
    Changing the lift changes the result by a boundary.
  \end{claim}

  \begin{case}[The class is zero]\label{case:section-zero}
    Choose a boundary representative and compute directly.
  \end{case}

  \begin{case}[The class is arbitrary]\label{case:section-general}
    Subtract two lifts and reduce to the first case.
  \end{case}
\end{proof}

References include \cref{mainthm:section}, \cref{thm:section},
\cref{lem:section}, \cref{claim:section}, \cref{step:section-lift}, and
\cref{case:section-zero}. Draft structural references such as
\cref{sec:section-gallery} and \cref{subsec:section-gallery} omit tags.

% =============================================================================
% Compact feature preview
% =============================================================================
\subsection[tag=0SE0203,label=subsec:section-features]{Mathematics and code}

Inline mathematics remains calm in prose:
\(\mathcal{L}\otimes\mathscr{F}\in D^b(X)\) and
\(\mathfrak{m}_x/\mathfrak{m}_x^2\cong\mathbb{A}^n_k\).
The same notation works in a labeled display:
\begin{equation}\label{eq:section-complex}
  \cdots \longrightarrow C_{n+1}
  \xrightarrow{d_{n+1}} C_n
  \xrightarrow{d_n} C_{n-1}
  \longrightarrow \cdots .
\end{equation}

\[
\begin{tikzcd}
  0 \arrow[r] & A_\bullet \arrow[r] \arrow[d]
    & B_\bullet \arrow[r] \arrow[d]
    & C_\bullet \arrow[r] \arrow[d] & 0 \\
  0 \arrow[r] & A'_\bullet \arrow[r]
    & B'_\bullet \arrow[r]
    & C'_\bullet \arrow[r] & 0
\end{tikzcd}
\]

\begin{lstlisting}[
  language=Python,
  caption={A tiny section-level listing},
  label=lst:section-example
]
def homology(cycles, boundaries):
    return cycles.quotient(boundaries)
\end{lstlisting}

See \cref{eq:section-complex} and \cref{lst:section-example}.

\end{document}
